
%(BEGIN_QUESTION)
% Copyright 2015, Tony R. Kuphaldt, released under the Creative Commons Attribution License (v 1.0)
% This means you may do almost anything with this work of mine, so long as you give me proper credit

\vskip 10pt

Identifiser områder i studiet der du ønsker å bli sterkere. Eksempler er teknisk lesing, kretsanalyse, bestemte oppgavetyper, tidsstyring og/eller ferdigheter i lab. Vis gjerne konkrete eksempler og ta disse opp med instruktør for målrettet forbedring. Som start kan du se temaoversikten på første side av arbeidsarket for neste mestringseksamen, samt ``General Values and Expectations'' tidlig i arbeidsarket.

\vskip 20pt

Identifiser deretter praktiske strategier du vil bruke for å styrke disse områdene. Eksempler er fokus på bestemte problemløsningstyper i lekser, målrettet trening på spesifikke emner og/eller samarbeid med labteam for mer trening på konkrete ferdigheter.

\vskip 20pt

\vskip 20pt \vbox{\hrule \hbox{\strut \vrule{} {\bf Forslag til sokratisk diskusjon} \vrule} \hrule}

\begin{itemize}
\item{} En nyttig strategi er å føre {\it journal} over hva du lærer i emnet. Utforsk hvordan daglig arbeid kan omformes til journal eller blogg for refleksjon og portefølje.
\item{} Hvor ligger øvingsoppgavene på {\it Socratic Instrumentation}-nettstedet?
\item{} Gå gjennom ``feedback questions'' for å finne spørsmål relatert til områder du vil styrke.
\end{itemize}

\underbar{file i00999}
%(END_QUESTION)





%(BEGIN_ANSWER)


%(END_ANSWER)





%(BEGIN_NOTES)

\noindent
{\bf Forberedelsesquiz:}

Nevn to tema fra kommende mestringseksamen, unntatt {\it instrument range calculations (Question \#4)}, {\it algebra (Question \#5)} eller {\it circuit fault analysis (Question \#6)}.

\vfil \eject

\noindent
{\bf Forberedelsesquiz:}

Nevn noen grunnleggende prinsipper (fysikk, kjemi, matematikk, logikk, elektriske kretser, prosesstyring osv.) du har brukt så langt, og kort hvordan du brukte dem for problemløsning/forståelse.

%INDEX% Metacognitive log, prompting students to reflect on principles and learning strategies
%INDEX% Reading Apprenticeship technique, metacognitive log

%(END_NOTES)

